\documentclass[11pt]{article}
\usepackage[margin=1in]{geometry}
\usepackage{amsmath,amssymb,amsthm,booktabs}
\usepackage[hidelinks]{hyperref}
\usepackage{xurl}
\setlength{\emergencystretch}{3em}
\newtheorem{theorem}{Theorem}
\newtheorem{lemma}[theorem]{Lemma}
\newtheorem{corollary}[theorem]{Corollary}
\newcommand{\C}{\mathbb C}
\newcommand{\F}{\mathbb F}
\newcommand{\GL}{\operatorname{GL}}
\newcommand{\SL}{\operatorname{SL}}
\newcommand{\PSL}{\operatorname{PSL}}
\newcommand{\rank}{\operatorname{rank}}
\newcommand{\ar}{\operatorname{ar}_{\C}}
\newcommand{\kfin}{\kappa_{\mathrm{fin}}}
\newcommand{\dfin}{d_{\mathrm{fin}}}
\title{Finite three-dimensional reflection lifts of involutory knot quotients}
\author{Written lifting argument with primary exceptional-group imports}
\date{17 November 2026}
\begin{document}
\maketitle

For a knot $K$, let $\kfin(K)$ be the maximum dimension of an
irreducible finite-image complex KCH representation with nontrivial
meridian eigenvalue. Let $\dfin(K)$ be the maximum meridian-class
generating rank among its finite complex images with meridian
displacement rank one. The maxima and their equality are established
in \cite[Corollary~\texttt{cor:finitecapacity}]{domination}.
Ordinary $U=1$ complex augmentation rank is denoted $\ar(K)$; its
longitude parameter is arbitrary nonzero.

\begin{theorem}\label{thm:lifts}
Let $K$ be a knot. Suppose $\pi_1(S^3\setminus K)$ has an epimorphism
onto $A_6$, or onto $\GL_3(\F_2)$, taking a meridian to a nonidentity
involution. Then it has an irreducible three-dimensional finite-image
complex representation with meridian eigenvalues $-1,1,1$.
The constructed full image is $G_{27}$ in the $A_6$ case and $G_{24}$
in the $\GL_3(\F_2)$ case, in their natural Shephard--Todd actions.
In both cases the meridian class has generating rank exactly three.
In particular,
\[
                    \kfin(K)\ge3,\qquad \dfin(K)\ge3.
\]
\end{theorem}

The construction uses a Wirtinger presentation of a knot group.
It does not assert that unique involution lifts extend to
homomorphisms from every group normally generated by one element.

\section{The precise primary group identifications}

We use the following rank-three irreducible complex reflection groups:
\[
\begin{array}{c|r|c|c}
 W&|W|&Z=W\cap\C^*I& W/Z\\\hline
 G_{24}&336&\{I,-I\}&\GL_3(\F_2)\\
 G_{27}&2160&\{\zeta I:\zeta^6=1\}&A_6 .
\end{array}
\]
These projective identifications are imported as group
identifications, not inferred from their orders.

Shephard--Todd \cite[Section~3, printed p.~278]{st54} identify the
projective order-$360$ group as Valentiner's group and explicitly
state its isomorphism with the alternating group of degree six.
They give the corresponding order-$2160$ reflection group, generated
by $45$ involutory reflections. Their Table~VII, printed p.~301,
identifies this natural rank-three irreducible group as number~27.

For number~24, the same section identifies the projective group as
Klein's simple group, with order-$336$ reflection cover. Boalch
\cite[printed p.~1]{boalch} explicitly identifies Klein's subgroup
of $\operatorname{PGL}_3(\C)$ with $\PSL_2(7)$ and its order-$336$
cover as generated by three involutory reflections.
The ATLAS page \cite{atlas} explicitly identifies $\PSL_2(7)$ with
$\PSL_3(2)$; the latter is $\GL_3(\F_2)$ because $\F_2^*$ and the
scalar center of $\SL_3(\F_2)$ are trivial.
Shephard--Todd's Table~VII row~24 identifies this cover with $G_{24}$.

The scalar-kernel orders follow from the displayed full and
projective orders. A finite subgroup of $\C^*$ of order $e$ is the
group of all $e$th roots of unity, which gives the literal scalar
kernels above. Both groups are generated by involutory reflections,
so their determinant images are $\{1,-1\}$.

\section{Unique involutions in the determinant-one cover}

Fix either group $W$ above, its scalar kernel $Z$, and $P=W/Z$.
Write $p:W\to P$ for projectivization and put
\[
                         H=\ker(\det:W\to\{1,-1\}).
\]
Since $-I\in Z$ and $\det(-I)=-1$ in dimension three,
$p|_H$ is onto: multiplying any lift by $-I$ changes its determinant
and preserves its projective image. Its kernel is
\[
 A=H\cap Z=\begin{cases}
 \{I\},&W=G_{24},\\
 \{\zeta I:\zeta^3=1\},&W=G_{27}.
 \end{cases}
\]
Thus $H\to P$ is respectively an isomorphism or a central extension
with kernel of order three. No non-splitting or Schur-multiplier
claim is needed.

\begin{lemma}\label{lem:unique}
Every nonidentity involution $q\in P$ has a unique involutory lift
$t(q)\in H$. For $h\in H$,
\[
                   t(p(h)qp(h)^{-1})=h\,t(q)\,h^{-1}.
\]
The eigenvalues of $t(q)$ are $1,-1,-1$.
\end{lemma}
\begin{proof}
Take $h\in H$ above $q$. Then $h^2\in A$. The square map is an
automorphism of the odd-order central cyclic group $A$, so there
is a unique $a\in A$ with $a^2=h^{-2}$. The element $ah$ is the
unique involution above $q$ in $H$. Conjugation preserves the
property of being such a lift and proves the displayed formula.
The lift is a nonidentity order-two matrix of determinant one in
dimension three. Its eigenvalues are therefore $1,-1,-1$.
\end{proof}

\begin{lemma}\label{lem:class}
All reflections of $W$ are involutions and form a single conjugacy
class, in bijection with the nonidentity involutions of $P$.
Their class-generating rank is three.
\end{lemma}
\begin{proof}
A finite-order reflection has order equal to the order of its
determinant. Since $\det W=\{1,-1\}$, all reflections have order two.
Each has determinant $-1$ and nontrivial projective image: neither
scalar involution $I,-I$ is a reflection in dimension three.
For every projective involution $q$, the unique reflection above
it is $-t(q)$. Indeed this is a reflection by the spectrum in
Lemma~\ref{lem:unique}; conversely the negative of any reflection
above $q$ is an involution in $H$ and hence is $t(q)$.

The nonidentity involutions of $A_6$ are its double transpositions
and form a single class. They are conjugate in $S_6$, and an odd
conjugator can be made even by multiplication by a transposition
on the two unused symbols, which centralizes a double transposition.
For $\GL_3(\F_2)$, an involution is $I+N$ with $N\ne0$ and $N^2=0$.
Its rank is one, since $\operatorname{im}N\subset\ker N$ in dimension
three. All such matrices have Jordan type $(2,1)$ and form one class.
Conjugation and the unique reflection lifts now show that the
reflections of $W$ form one class.

Both $G_{24}$ and $G_{27}$ have three reflection generators by
\cite[Table~VII and its complete exception footnote, printed p.~301]{st54}.
For the lower bound, if reflections $s_1,\ldots,s_k$ generate $W$,
then
\[
                  L=\sum_i\operatorname{im}(s_i-I)
\]
is a nonzero invariant subspace: $s_jx-x$ belongs to the $j$th
summand for every $x\in L$. Irreducibility gives $L=\C^3$, so
$k\ge3$. The common reflection class consequently has rank three.
\end{proof}

\section{The knot lift and full finite image}

\begin{proof}[Proof of Theorem~\ref{thm:lifts}]
Let $\rho:\pi_1(S^3\setminus K)\twoheadrightarrow P$ be the given
quotient. Choose positive meridians $m_i$ in a Wirtinger
presentation, and put $q_i=\rho(m_i)$. All $q_i$ are nonidentity
involutions because knot meridians are conjugate.
Assign $m_i\mapsto t(q_i)\in H$.
For each signed crossing relation
\[
                    m_b=m_o^{\epsilon}m_a m_o^{-\epsilon},
                       \qquad \epsilon\in\{1,-1\},
\]
the two sides after this assignment are involutions in $H$ above
the same element of $P$. Lemma~\ref{lem:unique} makes them equal.
Thus the assignment defines a homomorphism
$\widetilde\rho:\pi_1(S^3\setminus K)\to H$ with
$p\widetilde\rho=\rho$.

Let $\ell:\pi_1(S^3\setminus K)\to\mathbb Z$ be abelianization,
sending a positive meridian to one. The central scalar twist
\[
              \psi(g)=(-I)^{\ell(g)}\widetilde\rho(g)\in W
\]
is a representation. Every meridian has eigenvalues $-1,1,1$,
so its displacement rank is one. The image $J$ is finite, and
$p(J)=P$, since the twist is scalar.

In fact $J=W$. Projective surjectivity gives $W=JZ$.
For a meridian reflection $r\in J$ and any $w=jz\in JZ$,
\[
                       wrw^{-1}=jrj^{-1}\in J.
\]
Hence $J$ contains the entire $W$-class of $r$. By
Lemma~\ref{lem:class} this is every reflection of $W$, and those
reflections generate $W$. Thus $J=W$.
The natural action is irreducible of dimension three by the
primary classification, and the same lemma supplies class rank three.
\end{proof}

The lift $\widetilde\rho$ itself has meridian displacement rank two;
the scalar parity twist is necessary for the rank-one conclusion.
This construction provides finite images directly. It does not
require finite-image recognition of any earlier cyclotomic witness.

\section{The retained named knots and the sixteen exact capacities}

The exact knot-quotient inputs used here are already retained:
\path{ternary_affine_quotient_inputs.json} gives $24$ quotients whose
projection is $\GL_3(\F_2)$, and \path{a6_affine_inputs.json} gives
$24$ whose projection is $A_6$. Their projected meridians are
nonidentity involutions and the projections are onto; these are
parts of the positive quotient proofs
\cite[Theorem~\texttt{thm:knots}]{ternary} and
\cite[Theorems~\texttt{thm:sharp}, \texttt{thm:groups}]{a6}.
The two lists are disjoint. Theorem~\ref{thm:lifts} applies to all
$48$ named knots and proves $\kfin,\dfin\ge3$ for each.
No augmentation or capacity upper bound for all $48$ is asserted.

The existing generic obstruction covers precisely the following
sixteen members, with the indicated quotient types:
\[
\begin{array}{c|l}
 P&\text{names}\\\hline
 \GL_3(\F_2)&
 13a_{373},13a_{375},13a_{391},13a_{392},\\
 &13a_{3264},13a_{3266},13a_{3275},13a_{3276}\\[2pt]
 A_6&13a_{304},13a_{362},13a_{623},13a_{629},\\
 &13a_{3232},13a_{3259},13a_{3374},13a_{3375}.
\end{array}
\]
Their literal PD diagrams in the upper-bound input agree with those
in the corresponding quotient inputs.

\newpage
\begin{corollary}\label{cor:sixteen}
For every one of these sixteen knots,
\[
                 \dfin(K)=\kfin(K)=\ar(K)=3.
\]
In each case the finite class ceiling is attained by the constructed
$G_{24}$ or $G_{27}$ image at meridian eigenvalue $-1$.
\end{corollary}
\begin{proof}
Theorem~\ref{thm:lifts} gives a finite irreducible three-dimensional
KCH representation and a finite class bound of three.
The all-parameter irreducible-dimension upper bound and augmentation
upper bound for these exact sixteen diagrams are
\cite[Theorem~\texttt{thm:all16}, Corollary~\texttt{cor:augmentation}]{generic}.
Cornwell's KCH-to-augmentation construction and rank convention
\cite[Proposition~2.12, Definition~1.8]{cor17} give $\ar(K)\ge3$.
Finally finite complex reflection-class bounds are at most ordinary
complex augmentation rank by \cite[Theorem~\texttt{thm:domination}]{domination}.
All displayed bounds are therefore equal to three.
\end{proof}

The other $32$ quotient knots receive only the stated positive
finite-image lower bounds from this argument. No exact ordinary
augmentation rank is inferred for them. No finite-attainment
assertion for connected sums is made here.

\paragraph{AI assistance and review.}
ChatGPT (OpenAI) assisted with mathematical arguments, literature checks,
exposition, proof review and verification code.

\begin{thebibliography}{9}
\bibitem{st54} G.~C.~Shephard and J.~A.~Todd,
\emph{Finite unitary reflection groups}, Canad. J. Math.\ \textbf{6}
(1954), 274--304. \href{https://doi.org/10.4153/CJM-1954-028-3}{DOI}.
\bibitem{boalch} P.~Boalch, \emph{From Klein to Painlev\'e via Fourier,
Laplace and Jimbo}, \href{https://arxiv.org/abs/math/0308221v3}{arXiv:math/0308221v3}
(20 May 2004), printed p.~1.
\bibitem{atlas} ATLAS of Finite Group Representations,
\emph{Linear group $L_2(7)=L_3(2)$},
\url{https://brauer.maths.qmul.ac.uk/Atlas/lin/L32/}, retained 17 November 2026.
\bibitem{cor17} C.~R.~Cornwell, \emph{KCH representations, augmentations,
and A-polynomials}, J. Symplectic Geom.\ \textbf{15} (2017), 983--1017;
\href{https://arxiv.org/abs/1310.7526v3}{arXiv:1310.7526v3}.
\bibitem{ternary} Retained proof, \path{TERNARY_AFFINE_QUOTIENTS.tex}.
\bibitem{a6} Retained proof, \path{A6_AFFINE_SHARP_FOUR.tex}.
\bibitem{generic} Retained proof, \path{GENERIC_KCH_EXTENSIONS.tex}.
\bibitem{domination} Retained comparison,
\path{FINITE_COMPLEX_AUGMENTATION_DOMINATION.tex}.
\end{thebibliography}
\end{document}
